% Contenu partagé — les environnements à boîte et de type théorème.
%
% Ce fichier ne contient aucun préambule : il est inclus par plusieurs
% documents, qui l'affichent avec des thèmes différents. C'est la base de
% comparaison v1 / v2.

\lstset{language=[LaTeX]TeX}

\chapter{Boîtes et théorèmes}%
\label{chap:boxes}
\minitoc%

\begin{chapterintro}
    Les résultats (théorème, définition, proposition, corollaire, conjecture,
    lemme, théorème cité) partagent un compteur : deux résultats d'une même
    section portent donc des numéros distincts. Les exemples, les remarques et
    les exercices ont chacun leur propre compteur. Chaque famille a une variante
    étoilée, non numérotée. Toutes prennent une liste de clés optionnelle :
    `title=', `label=' et `note='. Une étiquette est posée telle quelle, sans
    préfixe ajouté par le template.
\end{chapterintro}

%---------------------------------------------------------------------------
\section{Les familles à cadre}%
\label{sec:framed}

\begin{lstlisting}
\begin{theorem}[title={Theoreme avec titre}, label=thm:cauchy]
    Un enonce. Cette boite porte l'etiquette brute \texttt{thm:cauchy} et se cite
    par \verb|\ref{thm:cauchy}| : theoreme~\ref{thm:cauchy}.
\end{theorem}
\end{lstlisting}
\begin{theorem}[title={Théorème avec titre}, label=thm:cauchy]
    Un énoncé. Cette boîte porte l'étiquette brute \texttt{thm:cauchy} et se cite
    par \verb|\ref{thm:cauchy}| : théorème~\ref{thm:cauchy}.
\end{theorem}

\begin{lstlisting}
\begin{theorem}
    Un enonce sans titre ni etiquette : la liste de cles est facultative.
\end{theorem}
\end{lstlisting}
\begin{theorem}
    Un énoncé sans titre ni étiquette : la liste de clés est facultative.
\end{theorem}

\begin{lstlisting}
\begin{definition}[title={Definition}, label=def:ouvert, note={utile}]
    Un enonce. Reference~: definition~\ref{def:ouvert}.
\end{definition}
\end{lstlisting}
\begin{definition}[title={Définition}, label=def:ouvert, note={utile}]
    Un énoncé. Référence~: définition~\ref{def:ouvert}.
\end{definition}

\begin{lstlisting}
\begin{proposition}
    Un enonce.
\end{proposition}
\end{lstlisting}
\begin{proposition}
    Un énoncé.
\end{proposition}

\begin{lstlisting}
\begin{corollary}[title={Corollaire}]
    Un enonce.
\end{corollary}
\end{lstlisting}
\begin{corollary}[title={Corollaire}]
    Un énoncé.
\end{corollary}

\begin{lstlisting}
\begin{conjecture}[title={Conjecture}]
    Un enonce. Cette famille venait du fork HDR et manquait au template de ce
    depot.
\end{conjecture}
\end{lstlisting}
\begin{conjecture}[title={Conjecture}]
    Un énoncé. Cette famille venait du fork HDR et manquait au template de ce
    dépôt.
\end{conjecture}

%---------------------------------------------------------------------------
\section{Les environnements de type théorème}%
\label{sec:theoremlike}

Ceux-là ne sont pas encadrés : ils se composent au fil du texte, intitulé en
gras suivi du texte sur la même ligne.

\begin{lstlisting}
\begin{lemma}
    Un lemme numerote dans la section.
\end{lemma}
\end{lstlisting}
\begin{lemma}
    Un lemme numéroté dans la section.
\end{lemma}

\begin{lstlisting}
\begin{lemma*}
    Un lemme non numerote --- variante etoilee.
\end{lemma*}
\end{lstlisting}
\begin{lemma*}
    Un lemme non numéroté --- variante étoilée.
\end{lemma*}

\begin{lstlisting}
\begin{example}
    Un exemple. Il se termine par un carre blanc.
\end{example}
\end{lstlisting}
\begin{example}
    Un exemple. Il se termine par un carré blanc.
\end{example}

\begin{lstlisting}
\begin{example}[title={Cas $E=\R$}]
    Un exemple dont le titre contient une macro mathematique --- garde-fou de
    la regression de l'issue \#18 : la detection \texttt{title=} forcait
    autrefois l'expansion complete du titre (\texttt{xstring}) et cassait des
    qu'une macro non trivialement expansible y apparaissait, meme standard
    (\texttt{\textbackslash R}, \verb|\mathbb{R}|...).
\end{example}
\end{lstlisting}
\begin{example}[title={Cas $E=\R$}]
    Un exemple dont le titre contient une macro mathématique --- garde-fou de
    la régression de l'issue \#18 : la détection \texttt{title=} forçait
    autrefois l'expansion complète du titre (\texttt{xstring}) et cassait dès
    qu'une macro non trivialement expansible y apparaissait, même standard
    (\texttt{\textbackslash R}, \verb|\mathbb{R}|...).
\end{example}

\begin{lstlisting}
\begin{example*}[note={avec une note}]
    Un exemple non numerote, avec une note entre parentheses.
\end{example*}
\end{lstlisting}
\begin{example*}[note={avec une note}]
    Un exemple non numéroté, avec une note entre parenthèses.
\end{example*}

\begin{lstlisting}
\begin{remark}
    Une remarque, marquee par un filet lateral.
\end{remark}
\end{lstlisting}
\begin{remark}
    Une remarque, marquée par un filet latéral.
\end{remark}

\begin{lstlisting}
\begin{remark*}
    Une remarque non numerotee.
\end{remark*}
\end{lstlisting}
\begin{remark*}
    Une remarque non numérotée.
\end{remark*}

\begin{lstlisting}
\begin{remark*}[Titre brut]
    Une variante etoilee avec un titre brut, sans `='. C'est la syntaxe
    encore en usage dans le corpus (par exemple \texttt{td/td3.tex}) : elle
    doit continuer a fonctionner, pas seulement la forme \texttt{title=}.
\end{remark*}
\end{lstlisting}
\begin{remark*}[Titre brut]
    Une variante étoilée avec un titre brut, sans `='. C'est la syntaxe
    encore en usage dans le corpus (par exemple \texttt{td/td3.tex}) : elle
    doit continuer à fonctionner, pas seulement la forme \texttt{title=}.
\end{remark*}

%---------------------------------------------------------------------------
\section{Compteurs séparés et variantes étoilées}%
\label{sec:counters}

Les résultats partagent un compteur, les exemples et les remarques ont chacun
le leur. Dans la suite ci-dessous, la définition et la proposition sont donc
numérotées 1 et 2, quel que soit le nombre d'exemples ou de remarques placés
entre elles : des diapositives qui reprennent ce polycopié sans la remarque
gardent exactement les mêmes numéros de résultats.

\begin{lstlisting}
\begin{definition}[title={Ensemble negligeable}]
    Un enonce.
\end{definition}
\begin{remark}
    Une remarque : son numero vient de son propre compteur.
\end{remark}
\begin{proposition}
    Un enonce, numerote juste apres la definition.
\end{proposition}
\end{lstlisting}
\begin{definition}[title={Ensemble négligeable}]
    Un énoncé.
\end{definition}
\begin{remark}
    Une remarque : son numéro vient de son propre compteur.
\end{remark}
\begin{proposition}
    Un énoncé, numéroté juste après la définition.
\end{proposition}

Chaque famille a une variante étoilée : même habillage, même titre, mais ni
numéro ni incrément de compteur. Elle sert à reprendre, dans des diapositives,
une remarque ou un exemple sans perturber les numéros du polycopié, ou à
rappeler un résultat d'un autre cours. Une étiquette posée sur une variante
étoilée n'aurait rien à désigner : elle est ignorée, avec un avertissement.

\begin{lstlisting}
\begin{theorem*}[title={Pythagore}]
    Un enonce rappele, non numerote.
\end{theorem*}
\end{lstlisting}
\begin{theorem*}[title={Pythagore}]
    Un énoncé rappelé, non numéroté.
\end{theorem*}

\begin{lstlisting}
\begin{definition*}
    Une definition non numerotee.
\end{definition*}
\end{lstlisting}
\begin{definition*}
    Une définition non numérotée.
\end{definition*}

\begin{proposition*}[note={rappel}]
    Une proposition non numérotée, avec une note.
\end{proposition*}

\begin{corollary*}
    Un corollaire non numéroté.
\end{corollary*}

\begin{conjecture*}
    Une conjecture non numérotée.
\end{conjecture*}

\begin{citedtheorem*}
    Un théorème cité, non numéroté.
\end{citedtheorem*}

Les numéros reprennent ensuite là où ils en étaient : la proposition suivante
suit directement celle du début de la section.

\begin{proposition}
    Un énoncé numéroté, qui suit la proposition précédente.
\end{proposition}

%---------------------------------------------------------------------------
\section{Un même énoncé, deux supports}%
\label{sec:portable}

C'est le point de la refonte. L'énoncé ci-dessous est copié tel quel depuis
l'exemple de TD~: en v0 il aurait fallu le réécrire, parce que
\texttt{theorem} et \texttt{question} n'avaient pas la même signature d'un
support à l'autre.

\begin{lstlisting}
\begin{exercise}[nosolution]
    Soit $E$ un espace vectoriel norme.
    \begin{question}
        Montrer que la boule $\ball(0,1)$ est convexe.
    \end{question}
    \begin{question}
        En deduire que $\norm{\cdot}$ est convexe.
    \end{question}
\end{exercise}
\end{lstlisting}
\begin{exercise}[nosolution]
    Soit $E$ un espace vectoriel normé.
    \begin{question}
        Montrer que la boule $\ball(0,1)$ est convexe.
    \end{question}
    \begin{question}
        En déduire que $\norm{\cdot}$ est convexe.
    \end{question}
\end{exercise}
